% TOP_PROBLEM: {"id": 20001103, "title": "Small sumsets with no three collinear", "category_id": 2, "category": {"id": 2, "name": "combinatorics", "display_name": "Combinatorics", "description": "Counting problems, graph theory, discrete structures.", "slug": "combinatorics", "order_index": 2, "created_at": "2026-07-31T15:26:25.671Z"}, "status": "partially_solved", "problem_number": "AIM-COMBINATORICS-0228", "set_id": 14}
% FIRST_POSTED: 2026-10-01
% ENTRY_KIND: statement_only
% UnsolvedMath ID 20001103; source catalogue code AIM-COMBINATORICS-0228
% !TeX program = lualatex
% Definition and sources only; this file is not an LLM solution attempt.
\documentclass[11pt,a4paper]{article}
\usepackage[margin=25mm]{geometry}
\usepackage{fontspec}
\setmainfont{lmroman10-regular.otf}[BoldFont=lmroman10-bold.otf,
 ItalicFont=lmroman10-italic.otf,BoldItalicFont=lmroman10-bolditalic.otf]
\usepackage{amsmath,amssymb,mathtools}
\usepackage{unicode-math}
\setmathfont{latinmodern-math.otf}
\AtBeginDocument{\let\setminus\smallsetminus}
\usepackage{xurl,hyperref}
\hypersetup{colorlinks=true,urlcolor=blue}
\setcounter{secnumdepth}{0}
\setlength{\parindent}{0pt}
\setlength{\parskip}{5pt}
\setlength{\emergencystretch}{3em}
\begin{document}
\section*{Small sumsets with no three collinear}\label{20001103}
{\small UnsolvedMath ID 20001103 · AIM-COMBINATORICS-0228 · Combinatorics · AIM Workshop Problem Lists}
\subsection{Definitions and mathematical statement}
For a finite set $A\subseteq\mathbb Z^2$, write
\[
                        A+A=\{a+b:a,b\in A\},
\]
where vector addition is coordinatewise, repeated summands are allowed,
and each resulting vector is counted once. Require that no affine
line in $\mathbb R^2$ contain three distinct points of $A$.
Equivalently, for distinct $a,b,c\in A$, the determinant
$\det(b-a,c-a)$ is nonzero. The condition forbids all collinear
triples, not merely equally spaced triples or triples on horizontal
or vertical lines.
For $n\ge1$, let
\[
 f_2(n)=\min\{|A+A|:A\subseteq\mathbb Z^2, |A|=n,
                                  A\text{ has no three collinear points}\}.
\]
Such sets exist for every $n$, for example by choosing integer points
on a parabola, and the minimum is over integer-valued cardinalities.

The problem asks for the smallest possible sumset size, including its
sharp asymptotic dependence on $n$. No bounding box, coordinate bound,
or convex-position assumption is present. Lower bounds must hold for
all eligible lattice configurations, while upper bounds require
configurations with many coinciding pair sums despite the prohibition
on collinearity.
For example, equality $a+b=c+d$ for different unordered pairs describes
a parallelogram relation; these additive coincidences need not produce
three collinear points. Consequently the geometric restriction does
not simply require every pair sum to be different. The target is the
minimum number of distinct sums, rather than the number of representations
or the area of the convex hull.
\subsection{Short English statement}
How small can the sumset of n lattice points be when no three of the points are collinear?
\subsection{Sources}
\begin{itemize}
\item[S1] ULAM.AI and the UnsolvedMath Contributors, supplied problems.json, record 20001103 (AIM-COMBINATORICS-0228), AIM Workshop Problem Lists. Catalogue identity and source transcription; CC BY 4.0.
\url{https://huggingface.co/datasets/ulamai/UnsolvedMath}.
\item[S2] American Institute of Mathematics, Recent trends in additive combinatorics, item 2.3. Original workshop question underlying AIM-COMBINATORICS-0228.
\url{https://aimath.org/WWN/additivecomb/additivecomb.pdf}.
\end{itemize}
\end{document}
